% Source printed pp. 512-513 / PDF pp. 520-521.
\subsection{\texorpdfstring{$G$-丛}{G-丛}}\label{src:10.2.3-subsection}
秩 $r$ 向量丛与主 $\mathrm{GL}_r$-丛之间有双射对应 (习题~\bookref{B.1.58}), 因而自然要将光滑曲线上多重稳定向量丛的模空间 $M_{r,d}^{\mathrm{ps}}(C)$ 推广到 $C$ 上的主 $G$-丛, 其中 $G$ 是任意约化群. 在代数闭域上, 主 $G$-丛的模叠 $\cBun_G(C)$ 是维数为 $(g-1)\dim G$ 的光滑代数叠, 其连通分支由 $\pi_1(G)$ 标记. 模叠 $\cBun_G(C)$ 在几何 Langlands 纲领中居于核心地位; 这一纲领将 $\cBun_G(C)$ 上 $\mathcal D$-模的导出范畴与 Langlands 对偶群的局部系统模叠上的拟凝聚层的导出范畴联系起来.

Ramanathan 将主 $G$-丛 $E$ 定义为半稳定的, 若对每个真抛物子群 $P\subseteq G$, 每个对应于截面 $\sigma:C\to E/P$ 的结构群到 $P$ 的约化, 以及每个支配特征 $\chi:P\to\mathbb G_m$, 都有 $\deg(\sigma^*L_\chi)\leq0$, 其中 $L_\chi$ 是 $E/P$ 上对应于 $\chi$ 的线丛. 这等价于要求存在忠实 $G$-表示 (相应地, 对每个 $G$-表示) $G\to\mathrm{GL}(V)$, 使对应的向量丛 $E(V)=E\times^G V$ 半稳定. 利用 GIT 可将半稳定主 $G$-丛的模空间构造为射影簇 \cite{Ram96}. 这一构造已被多次重新研究和推广 \cite{BS02,Sim94b,HS10}. 更详细的介绍见 \cite{Bal09,Sor99}.

\textbf{高维情形:} 对高维簇 $X$ 和一般的约化群 $G$, $G$-丛的模叠 $\cBun_G(X)$ 性质更差, 尚无已知的稳定性条件能从中截出射影模空间. 然而, 有若干方法利用无挠层扩充 $\cBun_G(X)$, 并依赖于忠实表示 $G\to\mathrm{GL}(V)$ \cite{Sch04,GS05,GLSS08}. 此时存在射影模空间, 但边界上的点不再参数化 $G$-丛.

\textbf{\texorpdfstring{$\overline M_g$}{Mgbar} 上的万有模空间:} 研究奇异稳定曲线上 $G$-丛的固有模空间, 以及构造 $\overline M_g$ 上的万有模空间, 是一个活跃的研究方向. 后者参数化对 $(C,E)$, 其中 $C$ 是稳定曲线, $E$ 是 $C$ 上的半稳定 $G$-丛. 目前已在 $G=\mathbb G_m$ \cite{Cap94} 和 $G=\mathrm{GL}_n$ \cite{Gie84,Pan96,NS99} 的情形实现.
